\section{Introdução}\label{sec:intro}

O cubo mágico é um dos exemplos mais conhecidos de \emph{quebra-cabeça de permutação}: seus estados são permutações de peças, seus movimentos geram um grupo finito e resolvê-lo equivale a encontrar uma palavra nos geradores que represente o inverso do estado atual. Essa formulação faz do cubo um laboratório natural para três áreas da Matemática Aplicada e Computacional: a \emph{teoria dos grupos}, que descreve quais estados são alcançáveis e como os movimentos se combinam~\cite{joyner2008,singmaster1981}; a \emph{teoria dos grafos}, que transforma o problema de resolução num problema de caminho mínimo num grafo de Cayley gigantesco~\cite{godsil2001,rokicki2014}; e a \emph{busca heurística}, que fornece algoritmos capazes de encontrar caminhos nesses grafos sem enumerá-los~\cite{hart1968,korf1985,korf1997}.

\paragraph{Motivação.} O projeto GraphCube~\cite{laguna2024} propôs uma visualização interativa elegante, em que famílias de círculos concêntricos representam as camadas do cubo e girar uma camada corresponde a deslizar pontos coloridos ao longo de um círculo. A proposta é visualmente atraente, mas não acompanha um modelo matemático explícito: não há representação de estados independente da tela, nem noção de grupo, distância ou algoritmo. Este trabalho parte dessa visualização para construir uma plataforma em que cada elemento gráfico tenha uma interpretação matemática precisa e verificável.

\paragraph{Trabalhos relacionados.} A estrutura de grupo do cubo é tratada em textos clássicos~\cite{singmaster1981,frey1982,joyner2008}. O diâmetro do grafo de Cayley do \cube{3} na métrica de meia-volta (HTM) foi determinado como 20 por Rokicki et al.~\cite{rokicki2014}, após uma série de cotas superiores como as 26 jogadas de Kunkle e Cooperman~\cite{kunkle2007}. Korf~\cite{korf1997} obteve as primeiras soluções ótimas para estados aleatórios combinando IDA$^*$~\cite{korf1985} com \textit{pattern databases}~\cite{culberson1998}, ideia depois estendida a heurísticas aditivas~\cite{felner2004}. O algoritmo de duas fases de Kociemba~\cite{kociemba1992,kociembaweb} encontra rapidamente soluções curtas, mas não necessariamente ótimas. Demaine et al.~\cite{demaine2011} estudaram a complexidade assintótica do cubo $n\times n\times n$, e abordagens recentes combinam aprendizado por reforço e busca~\cite{agostinelli2019}. Para o cálculo de ordens de grupos de permutação, o algoritmo de Schreier--Sims~\cite{sims1970,seress2003,holt2005} é a ferramenta padrão.

\paragraph{Objetivos e contribuições.} O objetivo não é produzir um resultado científico inédito sobre o cubo --- os resultados centrais são conhecidos --- mas construir uma ferramenta \emph{matematicamente rigorosa, reprodutível e didática} e documentar honestamente o que ela mostra. As contribuições deste trabalho são:
\begin{enumerate}[nosep,label=(\roman*)]
  \item uma análise do GraphCube que identifica precisamente o objeto que ele desenha (\cref{sec:graphcube}) e uma distinção formal entre grafo estrutural, grafo de estados e grafos de Cayley (\cref{sec:tres-grafos});
  \item um núcleo matemático em que todas as permutações derivam de uma descrição geométrica única, validado contra resultados da literatura (\cref{sec:modelagem,sec:res-validacao});
  \item a construção exata do diagrama de círculos do GraphCube como diagrama de incidência, generalizada para $n=2,3$ (\cref{sec:estrutural});
  \item implementações instrumentadas de BFS, BFS bidirecional, A$^*$, IDA$^*$ e do método de duas fases, com visualização que reproduz o registro real de cada execução (\cref{sec:desenvolvimento});
  \item quatro experimentos reprodutíveis sobre crescimento do espaço de busca, comparação de algoritmos, heurísticas e estrutura de grafos de Cayley (\cref{sec:metodologia,sec:resultados}).
\end{enumerate}

\paragraph{Procedência dos resultados.} Ao longo do texto, marcamos com \lit{} resultados conhecidos da literatura, com \ver{} resultados conhecidos que foram \emph{verificados computacionalmente} pelo nosso software e com \proj{} observações e medições produzidas especificamente neste projeto. Nenhuma afirmação de originalidade científica é feita além das contribuições de engenharia e de síntese acima.
